\documentclass[a4paper, 12pt]{article}

\usepackage{utils}

\renewcommand*{\today}{24 March 2026}

\begin{document}

\hotbox{Questions ouvertes}{CM 6}{\today}

\section{Annexes}

\subsection*{Annexe A: Code python générant des nombres tournants}

\begin{lstlisting}[breaklines=true]
def until_turned(first_value, mul, limit=100):
    value = [first_value]
    result = []
    retenu = 0
    
    for i in range(limit):
        if result and retenu == 0 and result[-1] == first_value:
            break
            
        prod = value[-1] * mul + retenu
        num = prod % 10
        retenu_tmp = prod // 10
        
        result.append(num)
        value.append(num)
        retenu = retenu_tmp
        
    original_str = ''.join(map(str, value[::-1][1:]))
    result_str = ''.join(map(str, result[::-1]))
    
    return original_str, len(value) - 1, result_str, len(result)

mul = 4
for i in range(1, 10):
    print(f"pour {i} avec mul = {mul} on a {until_turned(i, mul)}")
\end{lstlisting}

\subsection*{Annexe B: Caractérisation de la longueur d'un nombre $m$-tournant}

\begin{lemme}{}{existence}
    Soit $n \in \N^*$ et $p \in \Z/n\Z$.

    Si $p$ est inversible, alors $ord(p)$ existe.
\end{lemme}

\begin{demonstration}
    $\Z/n\Z$ est fini alors $\exists k, l \in \N, k > l$ tel que $p^k \equiv p^l [n]$.
    
    Ainsi, on a
    \begin{align*}        
        p^k \equiv p^l [n] \implies p^k (p^{-1})^l \equiv p^l (p^{-1})^l [n] \implies p^{k-l} \equiv 1 [n]
    \end{align*}

    Par conséquent, $ord(p) \leq k - l$
\end{demonstration}

\begin{theoreme}{Caractérisation de la longueur d'un nombre $m$-tournant}{}
    Soit $m \in \N^*$, $x$ le plus petit nombre $m$-tournant de chiffre des unités $d$ et $l$ sa longueur (nombre de chiffres).

    Alors $l = ord(10)$ dans $\Z/(\dfrac{10m - 1}{pgcd(d, 10m - 1)})\Z$.
\end{theoreme}

\begin{demonstration}
    On pose $y := \frac{x - d}{10}$ le nombre obtenu en retirant le chiffre des unités $d$ de $x$.

    Ainsi, on peut définir $x = 10y + d$.

    Par définition d'un nombre $m$-tournant, on a
    \begin{align*}
        &m \times x = d10^{l - 1} + y \\
        &\implies m(10y + d) = d10^{l - 1} + y \\
        &\implies m \times 10y - y = d10^{l - 1} - dm \\
        &\implies y(10m - 1) = d(10^{l - 1} - m) \\
        &\implies d(10^{l - 1} - m) \equiv 0 [10m - 1] \\
        &\implies d 10^{l - 1} \equiv dm [10m - 1]
    \end{align*}

    $d$ est inversible dans $\Z/(10m - 1)\Z$ si et seulement si \text{$pgcd(d, 10m - 1) = 1$}.
    Si $d$ n'est pas inversible, on peut poser $gcd = pgcd(d, 10m - 1)$ et on pourra considérer $d' = \dfrac{d}{gcd}$ et $m' = \dfrac{10m - 1}{gcd}$, ainsi $d'$ est inversible dans $\Z/m'\Z$.

    Ainsi, on a
    \begin{align*}
        d' 10^{l - 1} \equiv d'm [m'] &\implies 10^{l - 1} \equiv m [m'] \\
        &\implies 10^l \equiv 10m [m'] \\
        &\implies 10^l \equiv 1 [m']
    \end{align*}

    Donc $l = ord(10)$ dans $\Z/(\dfrac{10m - 1}{pgcd(d, 10m - 1)})\Z$.
\end{demonstration}

\subsection*{Annexe C: Existence d'un nombre tournant de toute longueur}

On se souvient de l'équation $y = \dfrac{d(m - 10^{l - 1})}{1 - 10m}$.

Si on pose $m = 10^{l - 1}$, alors $y = 0$ et $x = d$ est un nombre tournant de longueur $l$ pour tout $l \in \N^*$.

\subsection*{Annexe D: Base $b$ et rotation de $R$ quelconque}

\begin{theoreme}{}{}
    Soit $b, m, R \in \N^*$, $x$ le plus petit nombre $m$-$R$-tournant avec $d$ ses $R$ derniers chiffres et $l$ sa longueur.

    Alors $l = ord(b)$ dans $\Z/(\dfrac{b^R m - 1}{pgcd(d, b^R m - 1)})\Z$.
\end{theoreme}

\begin{demonstration}
    En reprenant la même démarche que pour une la base $10$ et une rotation de $1$, on pose $y := \dfrac{x - d}{b^R}$ le nombre obtenu en retirant les $R$ derniers chiffres de $x$.

    Ainsi, on a
    \begin{align*}
        m \times x = d b^{l - R} + y &\implies m(b^R y + d) = d b^{l - R} + y \\
        &\implies m \times b^R y - y = d b^{l - R} - dm \\
        &\implies y(b^R m - 1) = d(b^{l - R} - m) \\
        &\implies d(b^{l - R} - m) \equiv 0 [b^R m - 1] \\
        &\implies d b^{l - R} \equiv dm [b^R m - 1]
    \end{align*}

    $d$ est inversible dans $\Z/(b^R m - 1)\Z$ si et seulement si \text{$pgcd(d, b^R m - 1) = 1$}.
    Si $d$ n'est pas inversible, on peut poser $gcd = pgcd(d, b^R m - 1)$ et on pourra considérer $d' = \dfrac{d}{gcd}$ et $m' = \dfrac{b^R m - 1}{gcd}$, ainsi $d'$ est inversible dans $\Z/m'\Z$.

    Ainsi, on a
    \begin{align*}
        d' b^{l - R} \equiv d'm [m'] &\implies b^{l - R} \equiv m [m'] \\
        &\implies b^l \equiv b^R m [m'] \\
        &\implies b^l \equiv 1 [m']
    \end{align*}

    Donc $l = ord(b)$ dans $\Z/(\dfrac{b^R m - 1}{pgcd(d, b^R m - 1)})\Z$.
\end{demonstration}

\end{document}